\documentclass[10pt, a4paper]{article}

% ── PACKAGES ──────────────────────────────────────────────────
\usepackage[utf8]{inputenc}
\usepackage[T1]{fontenc}
\usepackage{geometry}
\geometry{
  top=2.0cm,
  bottom=2.0cm,
  left=2.2cm,
  right=2.2cm
}
\usepackage{hyperref}
\hypersetup{
  colorlinks=true,
  urlcolor=blue,
  linkcolor=blue,
  citecolor=blue,
  pdftitle={Pre-Registration P9: H3K27me3 Enrichment
            at Neural Crest Migration Loci in Domestic
            Sus scrofa vs. Wild Boar Embryos ---
            ChIP-seq Pre-Registration},
  pdfauthor={Eric Robert Lawson / OrganismCore},
  pdfsubject={Pre-registered experimental protocol
              and predictions},
  pdfkeywords={H3K27me3, EZH2, neural crest,
               domestication syndrome, ChIP-seq,
               Sus scrofa, wild boar, attractor geometry,
               developmental arrest, epigenetics,
               pre-registration, OrganismCore}
}
\usepackage{microtype}
\usepackage{parskip}
\usepackage{booktabs}
\usepackage{array}
\usepackage{tabularx}
\usepackage{longtable}
\usepackage{xcolor}
\usepackage{mdframed}
\usepackage{enumitem}
\usepackage{titlesec}
\usepackage{lmodern}
\usepackage{fancyhdr}
\usepackage{lastpage}
\usepackage{amsmath}
\usepackage{multirow}

% ── COLOURS ───────────────────────────────────────────────────
\definecolor{headergreen}{RGB}{20,100,60}
\definecolor{rulecolor}{RGB}{20,100,60}
\definecolor{claimbox}{RGB}{20,100,60}
\definecolor{notebox}{RGB}{240,248,240}
\definecolor{warningbox}{RGB}{255,248,220}
\definecolor{predbox}{RGB}{235,248,255}
\definecolor{scorebox}{RGB}{248,240,255}
\definecolor{failbox}{RGB}{255,235,235}
\definecolor{controlbox}{RGB}{240,255,240}
\definecolor{locusbox}{RGB}{245,240,255}
\definecolor{mechanismbox}{RGB}{248,240,255}

% ── HEADER / FOOTER ───────────────────────────────────────────
\pagestyle{fancy}
\fancyhf{}
\renewcommand{\headrulewidth}{0.4pt}
\fancyhead[L]{\small\color{headergreen}
  OrganismCore Pre-Registration P9 v1.0}
\fancyhead[R]{\small\color{headergreen}
  H3K27me3 / NCC Loci ChIP-seq --- 2026-03-18}
\fancyfoot[C]{\small Page \thepage\ of \pageref{LastPage}}

% ── SECTION FORMAT ────────────────────────────────────────────
\titleformat{\section}
  {\large\bfseries\color{headergreen}}
  {\thesection.}{0.5em}{}[
  \vspace{-0.3em}
  \textcolor{rulecolor}{\rule{\linewidth}{0.6pt}}]
\titleformat{\subsection}
  {\normalsize\bfseries\color{headergreen}}
  {\thesubsection.}{0.5em}{}
\titleformat{\subsubsection}
  {\small\bfseries}
  {\thesubsubsection.}{0.5em}{}

\setlength{\parskip}{0.4em}
\setlength{\parindent}{0pt}

% ══════════════════════════════════════════════════════════════
\begin{document}

% ── TITLE BLOCK ───────────────────────────────────────────────
\begin{center}
  {\LARGE\bfseries\color{headergreen}
  PRE-REGISTRATION P9}\\[0.4em]
  {\Large\bfseries
  H3K27me3 Enrichment at Neural Crest Migration
  Loci\\
  in Domestic \textit{Sus scrofa} vs.\
  Wild Boar Embryos:}\\[0.3em]
  {\large\bfseries
  The ChIP-seq Experiment}\\[0.4em]
  {\large
  EZH2-Mediated Epigenetic Silencing of the
  Neural Crest Migration Program\\
  as the Molecular Mechanism of
  Field-Specified Developmental Arrest\\
  Predicted and Locked Before Any
  Sequencing Data Exists}\\[0.6em]
  \noindent\textcolor{rulecolor}
  {\rule{\linewidth}{1pt}}\\[0.4em]
  {\normalsize
  \textbf{Author:} Eric Robert Lawson\\
  \textbf{Affiliation:} OrganismCore
  (Independent Research)\\
  \textbf{ORCID:}
  \href{https://orcid.org/0009-0002-0414-6544}
  {0009-0002-0414-6544}\\
  \textbf{Contact:}
  \href{mailto:OrganismCore@proton.me}
  {OrganismCore@proton.me}\\
  \textbf{Repository:}
  \href{https://github.com/Eric-Robert-Lawson/
  attractor-oncology}
  {\texttt{github.com/Eric-Robert-Lawson/
  attractor-oncology}}\\
  \textbf{Pre-registration timestamp:}
  2026-03-18\\
  \textbf{Theory paper:}
  \href{https://doi.org/10.5281/zenodo.19094935}
  {\texttt{doi:10.5281/zenodo.19094935}}\\
  \textbf{P7 pre-registration:}
  \href{https://doi.org/10.5281/zenodo.19096128}
  {\texttt{doi:10.5281/zenodo.19096128}}\\
  \textbf{Document version:} 1.0\\
  \textbf{License:}
  Creative Commons Attribution 4.0
  International (CC BY 4.0)}\\[0.4em]
  \noindent\textcolor{rulecolor}
  {\rule{\linewidth}{1pt}}
\end{center}

\vspace{0.5em}

% ── CLAIM BOX ─────────────────────────────────────────────────
\begin{mdframed}[
  backgroundcolor=claimbox,
  linecolor=claimbox,
  linewidth=0pt,
  innertopmargin=0.7em,
  innerbottommargin=0.7em,
  innerleftmargin=1.0em,
  innerrightmargin=1.0em
]
\color{white}
\textbf{THE PRE-REGISTERED CLAIM}

The molecular mechanism maintaining the
domestication developmental arrest is
EZH2-mediated H3K27me3 deposition at neural
crest migration driver loci during embryogenesis.
This is the same mechanism that maintains
cancer attractor arrest via EZH2-mediated
silencing of the luminal identity program in
triple-negative breast cancer.

\medskip
This pre-registration specifies: the loci to
be measured, the developmental stage at which
the measurement must be taken, the comparison
to be performed, the predicted direction of
enrichment, the scoring criteria, the control
conditions, and the success criteria ---
all locked before any ChIP-seq data is generated.

\medskip
\textbf{Primary prediction (P9):}
Domestic pig embryos at gestational day E14--E18
will show significantly elevated H3K27me3
enrichment at the following neural crest
migration driver loci relative to
wild boar embryos at the same developmental
stage with equivalent total H3K27me3 signal:
\textit{BMP4}, \textit{BMP7}, \textit{WNT1},
\textit{WNT3A}, \textit{CXCL12}, \textit{CXCR4},
\textit{EDNRB}, \textit{SOX10}, \textit{SNAI2}.
Minimum 4 of 9 target loci reach enrichment
ratio $\geq 1.5$ (domestic/wild boar)
at $q < 0.05$ (FDR-corrected).
\end{mdframed}

\vspace{1em}

% ══════════════════════════════════════════════════════════════
\section{Background and Theoretical Basis}
% ══════════════════════════════════════════════════════════════

\subsection{The Open Causal Question}

Wilkins, Wrangham and Fitch (2014) identified
that all domestication syndrome traits map to
a single developmental perturbation: mild
reduction in neural crest cell (NCC) migration
during embryogenesis. They attributed this
reduction to genetic change --- alleles affecting
NCC behaviour co-selected alongside tameness
alleles.

The field-arrest model
(\href{https://doi.org/10.5281/zenodo.19094935}
{Lawson, 2026, doi:10.5281/zenodo.19094935})
proposes that the causal origin is not genetic
but epigenetic: the domestication coherence
gradient field ($N_{\text{domestic}}$) suppresses
the maternal HPA axis, which alters the
intrauterine coherence gradient environment,
which establishes H3K27me3 enrichment at NCC
migration driver loci via EZH2 activity, which
arrests NCC migration at the juvenile-proximate
developmental stage.

\textbf{The molecular mechanism proposed is
EZH2-mediated H3K27me3 silencing of the NCC
migration program.} This is the same molecular
mechanism that maintains cancer attractor arrest
in triple-negative breast cancer (TNBC), where
EZH2 deposits H3K27me3 at \textit{FOXA1},
\textit{GATA3}, and \textit{ESR1} promoters,
silencing the luminal identity program and
maintaining the cancer attractor basin.

The primary field-level experimental prediction
of the field-arrest model is pre-registered
as P7
(\href{https://doi.org/10.5281/zenodo.19096128}
{doi:10.5281/zenodo.19096128}).
P9 is the molecular mechanism pre-registration:
it specifies the epigenetic signature that
must be present if the field-arrest mechanism
is correct, independently of whether P7
has yet been experimentally tested.

If P9 confirms, the molecular mechanism of
domestication and the molecular mechanism of
cancer are the same mechanism operating at
different scales:

\begin{center}
\small
\begin{tabular}{p{3.2cm}p{5.0cm}p{5.0cm}}
\toprule
& \textbf{Cancer (TNBC)} &
  \textbf{Domestication (\textit{Sus scrofa})} \\
\midrule
\textbf{Scale} &
  Single cell &
  Whole organism \\[0.2em]
\textbf{Silenced program} &
  Luminal identity
  (\textit{FOXA1}, \textit{GATA3},
  \textit{ESR1}) &
  NCC migration program
  (\textit{BMP4/7}, \textit{WNT1/3A},
  \textit{CXCL12/CXCR4}, \textit{EDNRB},
  \textit{SOX10}, \textit{SNAI2}) \\[0.2em]
\textbf{Silencing mechanism} &
  EZH2 $\rightarrow$ H3K27me3 &
  EZH2 $\rightarrow$ H3K27me3
  (predicted) \\[0.2em]
\textbf{Reversal agent} &
  Tazemetostat &
  Tazemetostat (predicted --- P11) \\[0.2em]
\textbf{Arrest phenotype} &
  Cancer &
  Domestic pig \\[0.2em]
\textbf{Resumed phenotype} &
  Luminal differentiation &
  Wild boar morphology \\
\bottomrule
\end{tabular}
\end{center}

\subsection{Why the NCC Migration Stage
E14--E18 Is the Critical Window}

Neural crest cell migration in \textit{Sus scrofa}
follows the standard vertebrate timeline:
NCC specification at the dorsal neural tube begins
at approximately E8--E10; delamination and active
migration at E10--E14; differentiation into
target tissues at E14--E20. The E14--E18 window
represents the period of maximum active migration,
when the gradient guidance environment (BMP, Wnt,
CXCL12/CXCR4) is most critical for migration
trajectory and therefore most sensitive to
epigenetic modification of the guidance genes.

H3K27me3 enrichment at migration driver loci at
E14--E18 would silence the gradient guidance
signals during their peak activity window,
producing partial migration arrest at the
juvenile-proximate stage without complete
ablation (which would produce lethal NCC deficiency
rather than viable domestic phenotype).

The E14--E18 window is therefore the predicted
critical measurement window. Earlier (E8--E12)
would capture specification but not guidance
arrest. Later (E20+) would capture
post-differentiation states rather than
the migration arrest itself.

\subsection{Why This Matters Beyond Domestication}

If H3K27me3 at NCC migration loci in domestic
embryos is confirmed, this establishes:

\begin{enumerate}[itemsep=0.2em]
  \item The molecular mechanism of domestication
    is epigenetic, not genetic. The genome was
    never modified. The epigenome was.

  \item The same enzyme (EZH2) that maintains
    cancer attractor arrest at the single-cell
    scale maintains domestication developmental
    arrest at the organismal scale. Scale-invariant
    epigenetic arrest mechanism confirmed.

  \item Tazemetostat --- currently in clinical use
    for EZH2-mutant cancers --- is the predicted
    reversal agent for domestication arrest (P11).
    P9 confirmation creates the biological
    justification for P11.

  \item The field ($N$) $\rightarrow$ epigenome
    $\rightarrow$ developmental trajectory
    causal chain is confirmed at the molecular
    level. The attractor geometry framework
    ($S + N + G \rightarrow R$) is grounded in
    a specific, testable molecular mechanism.
\end{enumerate}

% ══════════════════════════════════════════════════════════════
\section{The Pre-Registered Locus Set}
\label{sec:loci}
% ══════════════════════════════════════════════════════════════

\subsection{Primary Target Loci}

The following loci are pre-registered as the
primary measurement targets. They were selected
before literature consultation on \textit{Sus scrofa}
ChIP-seq data, based solely on their established
roles in vertebrate NCC migration guidance.

\begin{mdframed}[
  backgroundcolor=locusbox,
  linecolor=headergreen,
  linewidth=0.8pt,
  innertopmargin=0.5em,
  innerbottommargin=0.5em,
  innerleftmargin=0.8em,
  innerrightmargin=0.8em
]
\small
\begin{longtable}{p{2.0cm}p{4.5cm}p{6.0cm}}
\toprule
\textbf{Gene} &
\textbf{NCC Function} &
\textbf{Rationale for Inclusion} \\
\midrule
\textit{BMP4} &
NCC migration guidance gradient ---
dorsoventral axis &
BMP4 establishes the dorsoventral coherence
gradient that directs NCC migration away from
the neural tube. Silencing disrupts the
primary spatial guidance signal. \\[0.3em]
\textit{BMP7} &
NCC migration guidance gradient ---
complementary to BMP4 &
BMP7 acts as a complementary gradient signal
to BMP4 in the NCC migration environment.
Co-enrichment at both BMP4 and BMP7 would
indicate disruption of the complete BMP
guidance axis. \\[0.3em]
\textit{WNT1} &
NCC specification and
delamination trigger &
WNT1 is required for NCC specification and
initiates the delamination program. H3K27me3
enrichment here would arrest the program
before migration begins. \\[0.3em]
\textit{WNT3A} &
NCC specification amplification
and dorsal migration support &
WNT3A amplifies WNT1 signalling in
post-delamination NCC. Enrichment
here reduces migration competence. \\[0.3em]
\textit{CXCL12} &
Chemokine gradient --- directional
migration guidance &
CXCL12 (SDF-1) establishes the chemotactic
gradient that directs NCC to target tissues.
The CXCL12/CXCR4 axis is the primary
directional guidance mechanism.
Silencing of CXCL12 in source tissue
eliminates the gradient. \\[0.3em]
\textit{CXCR4} &
Chemokine receptor on migrating NCC ---
gradient sensor &
CXCR4 on NCC cells detects the CXCL12
gradient. H3K27me3 enrichment at CXCR4
reduces NCC responsiveness to the
directional signal. \\[0.3em]
\textit{EDNRB} &
Endothelin receptor B ---
melanocyte and enteric NCC migration &
EDNRB is specifically required for
melanocyte precursor and enteric nervous
system NCC migration. Enrichment here
specifically explains pigmentation
(depigmentation syndrome trait) and
connects directly to Marcstr\"{o}m's
white patch observation. \\[0.3em]
\textit{SOX10} &
NCC identity maintenance transcription
factor &
SOX10 maintains multipotency and migration
competence of NCC during migration.
H3K27me3 at SOX10 would cause NCC to
exit the migration-competent state
prematurely. \\[0.3em]
\textit{SNAI2} &
Epithelial-to-mesenchymal transition
(EMT) factor for NCC delamination &
SNAI2 (Slug) drives the EMT that allows
NCC to delaminate from the neural tube
and begin migration. H3K27me3 enrichment
here arrests delamination. \\
\bottomrule
\end{longtable}
\end{mdframed}

\subsection{Secondary Target Loci}

The following loci are pre-registered as
secondary targets. Enrichment here would
extend the prediction beyond NCC migration
to broader developmental programs.
They are not included in the primary
success criterion but are reported as
exploratory:

\begin{center}
\small
\begin{tabular}{p{2.2cm}p{10.5cm}}
\toprule
\textbf{Gene} & \textbf{Rationale} \\
\midrule
\textit{TWIST1} &
NCC EMT regulator, craniofacial development.
H3K27me3 here would explain snout shortening. \\[0.2em]
\textit{FOXD3} &
NCC multipotency maintenance.
Early NCC identity factor. \\[0.2em]
\textit{PAX3} &
NCC specification at the neural plate border.
Upstream of SOX10. \\[0.2em]
\textit{TFAP2A} &
NCC specification and survival during
migration. \\[0.2em]
\textit{DLX2} &
Craniofacial NCC differentiation.
H3K27me3 here would explain skull shape
changes. \\[0.2em]
\textit{PHOX2B} &
Autonomic nervous system NCC derivation.
Adrenal chromaffin cell specification. \\
\bottomrule
\end{tabular}
\end{center}

\subsection{Negative Control Loci}

The following loci are pre-registered as
negative controls. H3K27me3 enrichment is
\textit{not} predicted at these loci.
If enrichment is found at negative control
loci at comparable levels to positive
target loci, this indicates a global
increase in H3K27me3 rather than
locus-specific enrichment, which would
require re-analysis with stricter
normalisation.

\begin{center}
\small
\begin{tabular}{p{2.2cm}p{10.5cm}}
\toprule
\textbf{Gene} &
\textbf{Rationale for Selection as
Negative Control} \\
\midrule
\textit{GAPDH} &
Housekeeping gene. Should not be
H3K27me3 enriched in any active cell type. \\[0.2em]
\textit{ACTB} &
Housekeeping gene. Standard negative
control for histone modification ChIP. \\[0.2em]
\textit{MYF5} &
Skeletal muscle specification factor.
Not expressed in NCC or relevant to
NCC development. Should not be enriched. \\[0.2em]
\textit{HBB} &
Haemoglobin beta. Erythroid lineage
marker. Not expressed in embryonic
NCC tissue. \\
\bottomrule
\end{tabular}
\end{center}

% ══════════════════════════════════════════════════════════════
\section{Experimental Design}
\label{sec:design}
% ══════════════════════════════════════════════════════════════

\subsection{Overview}

H3K27me3 ChIP-seq on embryonic neural tube
and adjacent embryonic tissue (NCC migration
zone) from domestic \textit{Sus scrofa}
(commercial breed) and wild boar
(\textit{Sus scrofa} wild-caught, confirmed
no domestic admixture) at gestational day
E14--E18. Comparison: H3K27me3 enrichment
at pre-registered NCC migration driver loci,
domestic vs.\ wild, same developmental stage,
same tissue type.

No genetic modification. No
experimental intervention on the animals.
Pure comparison of epigenetic state at
equivalent developmental stage in
domestic vs.\ wild genome-matched tissue.

\subsection{Animal Subjects}

\textbf{Domestic embryos:}
Commercial breed pregnant sows
(Large White or Landrace, confirmed
breed registration). Embryos collected
at confirmed E14, E16, and E18
(timed from confirmed mating).
Minimum $n = 3$ litters per timepoint.
Minimum $n = 3$ embryos per litter
(one per litter used for ChIP;
remaining for RNA-seq validation ---
see Section~3.4).

\textbf{Wild boar embryos:}
Wild-caught pregnant \textit{Sus scrofa}
sows, confirmed by SNP panel to have
$< 5\%$ domestic admixture (same standard
as P7,
\href{https://doi.org/10.5281/zenodo.19096128}
{doi:10.5281/zenodo.19096128}).
Gestational age confirmed by
crown-rump length measurement against
established \textit{Sus scrofa} growth
curves. Embryos collected at equivalent
E14, E16, and E18. Minimum $n = 3$
litters per timepoint.

\textbf{Tissue collection:}
Embryos collected immediately post-euthanasia
of dam (licensed, IACUC-approved procedure).
Dissection within 30 minutes of collection.
Neural tube region (dorsal cranial neural tube
+ dorsal trunk neural tube, E14--E18 stage)
dissected on ice, snap-frozen in liquid
nitrogen within 60 minutes of collection.
Parallel tissue for RNA-seq fixed in RNAlater.

\subsection{ChIP-seq Protocol}

\subsubsection{Chromatin Preparation}

Tissue (neural tube + adjacent NCC migration
zone, approximately 50--100\,mg per sample)
cross-linked with 1\% formaldehyde, 10 minutes
at room temperature. Quenched with glycine.
Nuclei isolated by Dounce homogenisation
and differential centrifugation.
Chromatin sheared to 200--500\,bp by sonication
(Bioruptor Pico or equivalent: 30 cycles,
30\,s on/30\,s off at 4\textdegree C).
Shear size confirmed by Bioanalyzer before
proceeding.

\subsubsection{Immunoprecipitation}

Anti-H3K27me3 antibody:
Cell Signaling Technology \#9733
(rabbit monoclonal C36B11 --- the validated
antibody used across the domestication
syndrome field's epigenetic studies and
in all referenced cancer ChIP-seq work)
or Abcam ab192985 (rabbit monoclonal).
Both antibodies must be tested in parallel
if sufficient material allows.

Negative IP control: rabbit IgG isotype
control (same species and format as primary
antibody, same concentration).

Positive IP control: H3K27me3 at
\textit{CDKN2A} (p16/INK4a locus) ---
a validated PRC2/EZH2 target enriched by
H3K27me3 in multiple species and cell types.
Must confirm $>$ 10-fold enrichment over IgG
at \textit{CDKN2A} before any sample
is considered valid.

ChIP input: 5\% of chromatin retained
as input control per sample.
Library preparation: NEBNext Ultra II
ChIP-seq Library Prep Kit or equivalent.
Sequencing: Illumina paired-end 150\,bp,
minimum 30 million reads per sample.

\subsubsection{Sequencing Quality and
Alignment}

Quality control: FastQC on all raw reads
before alignment. Samples with $<$ 25 million
mappable reads or $>$ 30\% duplication rate
excluded and resequenced.

Alignment: BWA-MEM2 to the \textit{Sus scrofa}
reference genome (Sscrofa11.1, Ensembl release
109 or equivalent current release).
Duplicate removal: Picard MarkDuplicates.
Blacklist regions excluded using
ENCODE \textit{Sus scrofa} blacklist
(or ENCODE porcine blacklist if available;
otherwise mammalian consensus blacklist
lifted over to Sscrofa11.1).

\subsubsection{Peak Calling and
Differential Analysis}

Peak calling: MACS2 (--broad flag for
H3K27me3 broad domain marks,
$q$-value threshold 0.05).
Input subtracted per sample.

Differential enrichment at pre-registered
loci: DESeq2 on read counts in
$\pm$5\,kb windows centred on the
transcription start site (TSS) of each
target gene, plus any MACS2-called
peaks within 50\,kb of TSS.

Comparison: domestic (all timepoints) vs.\
wild boar (all timepoints), with gestational
day as covariate. Primary comparison:
E16 (mid-window) domestic vs.\ E16 wild boar.

FDR correction: Benjamini-Hochberg across
all pre-registered loci in a single
correction set. $q < 0.05$ threshold.
Enrichment ratio threshold: $\geq 1.5$
(domestic/wild boar) at the TSS$\pm$5\,kb
window.

\subsection{RNA-seq Validation}

For each embryo where ChIP material is
collected, parallel RNA-seq is performed
on adjacent neural tube tissue from the
same dissection.

\textbf{Rationale:} H3K27me3 enrichment at
a locus should correspond to reduced
transcription of that gene. RNA-seq confirms
that the predicted H3K27me3 enrichment
produces the predicted transcriptional
silencing. If H3K27me3 is elevated at
a locus but transcription is unchanged,
the enrichment is not functionally relevant.

\textbf{Protocol:}
Total RNA extraction: RNeasy Micro Kit
(Qiagen) or equivalent.
RNA quality: RIN $\geq$ 7.0 required
(Bioanalyzer).
Library: NEBNext Ultra II RNA-seq with
poly-A selection.
Sequencing: Illumina paired-end 150\,bp,
minimum 20 million reads per sample.
Alignment: STAR to Sscrofa11.1.
Differential expression: DESeq2.

\textbf{Predicted direction:}
Genes with elevated H3K27me3 in domestic
embryos (from ChIP-seq primary analysis)
should show \textit{reduced} expression
in domestic embryo neural tube relative
to wild boar at equivalent stage.
Confirmed H3K27me3 enrichment + reduced
expression = dual confirmation of silencing.

% ══════════════════════════════════════════════════════════════
\section{The Two Models and Their
Diverging Predictions}
\label{sec:models}
% ══════════════════════════════════════════════════════════════

\begin{center}
\small
\begin{tabular}{p{4.5cm}p{4.5cm}p{4.5cm}}
\toprule
\textbf{Outcome} &
\textbf{Genetic model} &
\textbf{Field-arrest model} \\
\midrule
H3K27me3 at NCC loci in
domestic vs.\ wild embryos
(same developmental stage) &
No difference predicted.
Genetic differences affect
allele-encoded programs,
not H3K27me3 patterns at
equivalent developmental stages. &
\textbf{Elevated in domestic embryos}
at $\geq$ 4/9 pre-registered loci.
Enrichment ratio $\geq 1.5$ at
$q < 0.05$. \\[0.4em]
EZH2 protein level in domestic
vs.\ wild embryo NCC zone &
No specific prediction. &
\textbf{Elevated in domestic embryos}
at E14--E18 in the NCC migration zone.
(Exploratory; not in primary success
criterion.) \\[0.4em]
Expression of NCC migration
drivers in domestic vs.\ wild
(RNA-seq) &
No difference predicted at
equivalent developmental stages. &
\textbf{Reduced in domestic embryos}
at loci with elevated H3K27me3. \\[0.4em]
H3K27me3 at negative control loci &
No enrichment predicted. &
No enrichment predicted.
(Same prediction --- used only
to detect global artefacts.) \\
\bottomrule
\end{tabular}
\end{center}

\vspace{0.3em}
\textbf{The distinguishing comparison is:}
H3K27me3 enrichment at NCC migration driver
loci in domestic vs.\ wild embryos at the
same developmental stage with equivalent
global H3K27me3 signal. This comparison
is not confounded by developmental stage
differences (same stage), by tissue
differences (same dissection protocol),
or by global chromatin accessibility
differences (input-subtracted ChIP,
housekeeping gene negative controls).

% ══════════════════════════════════════════════════════════════
\section{Success Criteria}
\label{sec:success}
% ══════════════════════════════════════════════════════════════

\subsection{Primary Success Criterion}

\begin{mdframed}[
  backgroundcolor=notebox,
  linecolor=headergreen,
  linewidth=1pt,
  innertopmargin=0.5em,
  innerbottommargin=0.5em,
  innerleftmargin=0.8em,
  innerrightmargin=0.8em
]
\textbf{P9-SUCCESS requires ALL of the
following:}

\medskip
\textbf{Criterion 1 (ChIP quality):}
$>$ 10-fold H3K27me3 enrichment over IgG
at the \textit{CDKN2A} positive control
locus confirmed in $\geq$ 80\% of samples
before any target locus analysis proceeds.

\medskip
\textbf{Criterion 2 (NCC locus enrichment):}
$\geq$ 4 of the 9 pre-registered primary
target loci show H3K27me3 enrichment ratio
(domestic/wild boar) $\geq 1.5$ at
$q < 0.05$ (BH-FDR corrected across all
9 loci) at the E16 primary timepoint.

\medskip
\textbf{Criterion 3 (Direction specificity):}
Mean H3K27me3 enrichment at negative
control loci (GAPDH, ACTB, MYF5, HBB)
does not differ significantly between
domestic and wild boar samples
($p > 0.05$, confirming the enrichment
at target loci is locus-specific, not
a global H3K27me3 increase).

\medskip
\textbf{Criterion 4 (RNA-seq concordance):}
$\geq$ 3 of the loci meeting Criterion 2
also show significantly reduced mRNA
expression in domestic embryo neural
tube relative to wild boar at the same
stage ($q < 0.05$, DESeq2),
confirming that H3K27me3 enrichment
corresponds to functional silencing.
\end{mdframed}

\subsection{Graded Confirmation Levels}

\begin{center}
\small
\begin{tabular}{p{3.5cm}p{9.5cm}}
\toprule
\textbf{Outcome} & \textbf{Interpretation} \\
\midrule
All 4 criteria met &
\textbf{Full confirmation.} EZH2/H3K27me3
silencing of the NCC migration program is
the molecular mechanism of domestication
developmental arrest. Proceed immediately
to P11 (tazemetostat experiment).
Publish P9 result and update theory paper
with molecular mechanism confirmed. \\[0.4em]
Criteria 1--3 met,
Criterion 4 not met &
\textbf{Structural confirmation, functional
confirmation absent.} H3K27me3 enrichment is
present at NCC loci but not reducing
transcription. Possible explanation: the
enrichment affects enhancer elements not
captured by TSS$\pm$5\,kb window. Extend
to enhancer-focused analysis (H3K27ac
ChIP-seq to map active enhancers in wild
boar; confirm their suppression in domestic).
Still considered partial confirmation. \\[0.4em]
Criterion 2 met (3/9 loci only,
below threshold) &
\textbf{Trend toward confirmation, below
pre-specified threshold.} Increase sample
size to $n \geq 5$ litters per timepoint
and repeat. If 3/9 becomes 4/9 at larger
$n$: full analysis. If remains $< 4/9$:
failure diagnosis. \\[0.4em]
Criterion 1 met,
Criterion 2 not met &
\textbf{True negative.} ChIP quality
confirmed (positive control passed).
H3K27me3 enrichment at NCC loci is absent.
EZH2/H3K27me3 is not the mechanism.
Failure diagnosis applied. \\
\bottomrule
\end{tabular}
\end{center}

% ══════════════════════════════════════════════════════════════
\section{Control Conditions and
Confound Management}
\label{sec:controls}
% ══════════════════════════════════════════════════════════════

\subsection{Critical Controls}

\begin{mdframed}[
  backgroundcolor=controlbox,
  linecolor=headergreen,
  linewidth=0.8pt,
  innertopmargin=0.5em,
  innerbottommargin=0.5em,
  innerleftmargin=0.8em,
  innerrightmargin=0.8em
]
\textbf{C1 --- Developmental stage matching
(primary confound):}
H3K27me3 patterns change substantially
across developmental stages. Domestic and
wild boar embryos must be confirmed at
equivalent developmental stage by
\textit{independent} criteria beyond gestational
day: crown-rump length measurement against
published \textit{Sus scrofa} staging tables,
and somite count where available.
Any sample pair with $>$ 1 day estimated
developmental age difference is excluded
from primary analysis and reported separately.

\medskip
\textbf{C2 --- Tissue dissection equivalence:}
The NCC migration zone is a small, specific
region of the embryo. Contamination by
non-NCC tissue (e.g., epidermis, notochord,
somites) will dilute the signal.
Dissection must be performed by the same
trained dissector for all samples.
Dissection protocol must be validated on
pilot samples by H\&E histology of adjacent
sections to confirm correct tissue composition.

\medskip
\textbf{C3 --- Antibody specificity:}
Two independent anti-H3K27me3 antibodies
must be used in parallel (CST \#9733 and
Abcam ab192985). Loci called as enriched
must show enrichment with both antibodies.
Any locus enriched with only one antibody
is reported as exploratory only.

\medskip
\textbf{C4 --- Input normalisation:}
5\% input control retained per sample.
All enrichment ratios computed after
input subtraction. Samples with input
library quality $<$ 15 million mappable
reads excluded.

\medskip
\textbf{C5 --- Domestic admixture exclusion:}
Wild boar samples must be confirmed at
$< 5\%$ domestic admixture by SNP panel
(same standard as P7,
\href{https://doi.org/10.5281/zenodo.19096128}
{doi:10.5281/zenodo.19096128}).
Any wild boar sample with $\geq 5\%$
admixture excluded. This controls for
the possibility that wild boar samples
with partial domestic ancestry show
partial H3K27me3 enrichment that
contaminates the wild boar baseline.

\medskip
\textbf{C6 --- Cross-laboratory replication:}
Primary analysis performed at one
sequencing facility. Results submitted
for replication at a second independent
facility on a minimum of 3 embryo pairs
(1 domestic, 1 wild boar, matched stage)
before primary success criterion is declared.
\end{mdframed}

\subsection{Pre-specified Confound
Interpretations}

\textbf{Confound: Breed-specific H3K27me3
differences unrelated to NCC.}
If all analysed loci (including negative
controls) show elevated H3K27me3 in domestic
vs.\ wild, this is a global methylation
difference, not a locus-specific NCC arrest
signal. The negative control criterion
(Criterion 3) specifically detects this.
If Criterion 3 fails, the dataset is invalid
for the primary claim and requires
reanalysis with global H3K27me3
normalisation (spike-in normalisation
with Drosophila chromatin).

\textbf{Confound: Different NCC developmental
timing between domestic and wild embryos.}
If domestic pigs have an accelerated or
delayed NCC migration timeline relative
to wild boar (independent of the arrest),
then E16 domestic and E16 wild boar
may not be at equivalent NCC migration
stages. This is detected by RNA-seq:
the expression signature of NCC migration
stage markers (\textit{SOX10}, \textit{SNAI2},
\textit{FOXD3} expression in the neural tube)
must be confirmed equivalent between
domestic and wild boar samples included
in the primary comparison.

% ══════════════════════════════════════════════════════════════
\section{Failure Diagnosis}
\label{sec:failure}
% ══════════════════════════════════════════════════════════════

\begin{mdframed}[
  backgroundcolor=failbox,
  linecolor=red!50!black,
  linewidth=0.8pt,
  innertopmargin=0.5em,
  innerbottommargin=0.5em,
  innerleftmargin=0.8em,
  innerrightmargin=0.8em
]
\textbf{If P9 fails (Criterion 1 passes,
Criterion 2 does not), the following
diagnoses are pre-specified:}

\medskip
\textbf{Diagnosis F1 (Wrong histone mark):}
The arrest is maintained by a different
Polycomb mechanism. H3K27me3 is the
canonical PRC2 mark but PRC1 (H2AK119ub)
or DNA methylation (5mCpG) could
alternatively silence the NCC migration
program.
\textit{Direction:} Perform H2AK119ub
ChIP-seq and WGBS on the same tissue pairs.
If H2AK119ub is enriched at NCC loci in
domestic vs.\ wild boar, the mechanism is
PRC1 rather than PRC2. The field-arrest
model is confirmed; the specific enzyme
is PRC1 rather than EZH2.

\medskip
\textbf{Diagnosis F2 (Wrong tissue:
arrest occurs earlier):}
The epigenetic arrest is established at
a pre-neural-crest stage (E8--E10,
specification stage) rather than during
active migration (E14--E18). H3K27me3
at migration guidance loci is secondary
to the primary arrest elsewhere.
\textit{Direction:} Perform the same
ChIP-seq protocol on E10--E12 embryos
(specification stage) to test for
earlier enrichment.

\medskip
\textbf{Diagnosis F3 (Wrong loci):}
The H3K27me3 arrest targets a different
set of NCC-related genes not in the
pre-registered panel.
\textit{Direction:} Perform genome-wide
differential H3K27me3 analysis (not
restricted to pre-registered loci).
Identify which, if any, NCC-related loci
show enrichment. Report as exploratory.

\medskip
\textbf{Diagnosis F4 (Mechanism is not
epigenetic):}
The NCC migration reduction in domestic
animals is genetic (allele-encoded),
not epigenetic. The field does not
establish H3K27me3. The genome specifies
reduced NCC migration capacity directly.
\textit{Direction:} Perform whole-genome
sequencing of domestic vs.\ wild boar
embryos. Identify coding variants in
the pre-registered NCC loci. If
functional variants in NCC migration
genes are present and track with
domestication, the genetic model is
supported over the field-arrest model
for this specific mechanism.
This outcome would be published as
a partial falsification of the
epigenetic arm of the field-arrest model
while leaving the basin-geometry
framework intact.
\end{mdframed}

% ══════════════════════════════════════════════════════════════
\section{Pathway to P11}
\label{sec:p11}
% ══════════════════════════════════════════════════════════════

P9 is the molecular prerequisite for P11
(tazemetostat in pregnant domestic sow,
shifting F1 morphology toward wild-type).
The logical dependency is:

\vspace{0.3em}
\begin{center}
\small
\begin{tabular}{p{4.5cm}p{0.5cm}p{8.0cm}}
\textbf{P9 confirms H3K27me3 at NCC loci
in domestic embryos} &
$\Rightarrow$ &
\textbf{EZH2 is the arrest maintenance
enzyme.}
Tazemetostat (EZH2 inhibitor) is the
predicted reversal agent.
P11 is justified and ready to proceed. \\[0.5em]
\textbf{P9 fails (Diagnosis F1:
PRC1/H2AK119ub instead)} &
$\Rightarrow$ &
\textbf{Tazemetostat is the wrong drug.}
PRC1 inhibitor (e.g., PTC-596 targeting
BMI1) is the predicted reversal agent.
P11 is redesigned with PRC1 inhibitor. \\[0.5em]
\textbf{P9 fails (Diagnosis F4:
genetic mechanism)} &
$\Rightarrow$ &
\textbf{P11 is not justified.}
No pharmacological reversal of the
arrest is predicted. The field-arrest
epigenetic model is falsified for
this specific mechanism.
\end{tabular}
\end{center}

\vspace{0.3em}
P9 is therefore not only a test of the
field-arrest mechanism. It is the gateway
experiment that determines the molecular
target for P11 and, by extension, provides
the biological justification for the
cross-scale connection between domestication
biology and EZH2 inhibitor cancer therapy.

% ══════════════════════════��═══════════════════════════════════
\section{Cost and Timeline}
\label{sec:cost}
% ══════════════════════════════════════════════════════════════

\begin{center}
\small
\begin{tabular}{p{7.0cm}p{5.0cm}}
\toprule
\textbf{Item} & \textbf{Estimated Cost (USD)} \\
\midrule
Wild boar embryo collection
(licensed, $n = 9$ litters, 3 per timepoint) &
\$8,000--\$18,000 \\[0.2em]
Domestic pig embryo collection
($n = 9$ litters, abattoir or farm) &
\$3,000--\$6,000 \\[0.2em]
Chromatin preparation and ChIP
(2 antibodies $\times$ 2 groups
$\times$ 3 timepoints $\times$ 3 replicates) &
\$12,000--\$20,000 \\[0.2em]
ChIP-seq library prep and sequencing
(36 samples $\times$ 30M reads) &
\$18,000--\$28,000 \\[0.2em]
RNA-seq library prep and sequencing
(36 samples $\times$ 20M reads) &
\$10,000--\$16,000 \\[0.2em]
Bioinformatic analysis
(alignment, peak calling, DESeq2) &
\$4,000--\$8,000 \\[0.2em]
SNP genotyping for admixture confirmation &
\$2,000--\$4,000 \\[0.2em]
Cross-laboratory replication
(3 embryo pairs) &
\$8,000--\$12,000 \\[0.2em]
Histological validation of tissue dissection &
\$2,000--\$4,000 \\[0.2em]
\midrule
\textbf{Total estimated cost} &
\textbf{\$67,000--\$116,000} \\
\bottomrule
\end{tabular}
\end{center}

\vspace{0.3em}
\textbf{Timeline from ethical approval
to primary data:}
3--6 months (embryo collection is the
rate-limiting step; sequencing and
analysis: 6--8 weeks after collection).

\textbf{Minimum timeline to replication:}
8--10 months.

% ═══════��══════════════════════════════════════════════════════
\section{Reporting Standard}
\label{sec:reporting}
% ══════════════════════════════════════════════════════════════

Results of this experiment will be reported
against this pre-registration regardless of
outcome. The following standards apply:

\begin{enumerate}[itemsep=0.2em]
  \item All 9 pre-registered primary target
    loci are reported. No selective reporting
    of only confirming loci.

  \item All 4 negative control loci are
    reported. If any negative control locus
    shows enrichment at threshold, this is
    reported prominently, not buried.

  \item The analysis is performed exactly as
    specified in Section~3.3. Any deviation
    must be declared in the methods with
    justification.

  \item The cross-laboratory replication
    result must be reported in the same
    publication as the primary result, not
    deferred to a separate publication.

  \item If P9 fails, the failure diagnosis
    and the follow-on direction are published
    in the primary report. A failure report
    is a result. It will be published.

  \item Sequencing data deposited in GEO
    (Gene Expression Omnibus) before
    submission to journal, with accession
    number included in the methods.
\end{enumerate}

% ══════════════════════════════════════════════════════════════
\section{Linked Records}
\label{sec:linked}
% ══════════════════════════════════════════════════════════════

The following records existed at the time
of this deposit and are cited with confirmed
DOIs. No forward references to undeposited
documents are included.

\begin{center}
\small
\begin{tabular}{p{2.8cm}p{4.2cm}p{5.0cm}}
\toprule
\textbf{ID} & \textbf{Title} & \textbf{DOI} \\
\midrule
Theory paper &
Domestication as Field-Specified
Developmental Arrest &
\href{https://doi.org/10.5281/zenodo.19094935}
{\texttt{10.5281/zenodo.19094935}} \\[0.4em]
P7 &
Cross-Fostering / NCC Arrest
Pre-Registration &
\href{https://doi.org/10.5281/zenodo.19096128}
{\texttt{10.5281/zenodo.19096128}} \\[0.4em]
P-INVERSION-1 &
Plant Architecture Inversion via
Coherence Gradient Field Specification &
\href{https://doi.org/10.5281/zenodo.19040399}
{\texttt{10.5281/zenodo.19040399}} \\[0.4em]
LECA-A &
LECA Resurrection Fungal Arm &
\href{https://doi.org/10.5281/zenodo.18986790}
{\texttt{10.5281/zenodo.18986790}} \\[0.4em]
FOXA1/EZH2 &
FOXA1/EZH2 Ratio as Universal
Breast Cancer Subtype Classifier &
\href{https://doi.org/10.5281/zenodo.18883922}
{\texttt{10.5281/zenodo.18883922}} \\
\bottomrule
\end{tabular}
\end{center}

% ── DOCUMENT METADATA ────────────────────────────────────────
\vspace{1em}
\noindent\textcolor{rulecolor}
{\rule{\linewidth}{0.8pt}}
\vspace{0.3em}

{\footnotesize
\textbf{Document ID:}
P9\_H3K27ME3\_NCC\_CHIPSEQ\_PRE\_
REGISTRATION\_v1.0\\
\textbf{Date:} 2026-03-18\\
\textbf{Author:} Eric Robert Lawson /
OrganismCore\\
\textbf{Theory paper:}
\href{https://doi.org/10.5281/zenodo.19094935}
{\texttt{doi:10.5281/zenodo.19094935}}\\
\textbf{P7 pre-registration:}
\href{https://doi.org/10.5281/zenodo.19096128}
{\texttt{doi:10.5281/zenodo.19096128}}\\
\textbf{Repository:}
\href{https://github.com/Eric-Robert-Lawson/
attractor-oncology}
{\texttt{github.com/Eric-Robert-Lawson/
attractor-oncology}}\\
\textbf{Document version:} 1.0 ---
LOCKED AT PRE-REGISTRATION.
No modifications after Zenodo deposit.\\
\textbf{License:} CC BY 4.0\\
\textbf{ORCID:}
\href{https://orcid.org/0009-0002-0414-6544}
{0009-0002-0414-6544}\\
\textbf{Contact:}
\href{mailto:OrganismCore@proton.me}
{OrganismCore@proton.me}
}

\end{document}
